% Physical unit 031; accepted Session 03 after fresh 400-dpi authority replay.
\GordanUnitStart{031}{printed page 26}
\GordanTopFolio{26}
\begingroup\small
Although Jacobi showed that this equation can be made to assume a far more elegant form, for our purpose we can nevertheless stand by it.

\sect{5}

Since the $\Theta$ functions are expanded into very rapidly convergent series, and therefore are computed very easily, it will now be useful also to refer the trigonometric functions of the angles
\[
 B\qquad l\qquad T
\]
to $\Theta$ functions.

The most important of the formulae developed in this way are these:
\[
\begin{gathered}
 \sin l=\sin l_0\sqrt{x'\over x}\,{H(u+K)\over\Theta(u)},\\
 \sqrt{x}\sin B=\sin B_0{H(u+K)\over\Theta(u+K)},\\
 {\cos B\over\cos l}=\sqrt{x'}{\Theta u\over\Theta(u+K)},\\
 \tg T=\cot l_0\sqrt{x'}\,{\Theta u\over H(u)}.
\end{gathered}
\]

\rule{0.35\textwidth}{0.4pt}

\caput{VII}
\sect{1}

But in order that the quantities
\[
 l\qquad s\qquad \psi
\]
may be brought to calculation from the equations of the preceding chapter, it is necessary to establish formulae for the constants
\[
 l_0\qquad B_0\qquad \Sigma(l_0)\qquad \mathcal A(l_0)\qquad x\qquad x'\qquad K\qquad E\qquad q\qquad v,
\]
all of them to be computed from one of them.

\sect{2}

Usually the circle that is touched by the curve is defined; only very rarely does one depart from this method. When one of the two periods is known, for example when it is asked: \gq{What circle is touched by the shortest line of all defined by two points of the equator?} then, for computing $l_0$, we use only the series of Chapter III.

\sect{3}

Likewise, now supposing $l_0$ known, we can determine the quantities
\[
 B_0\qquad \Sigma(l_0)\qquad \mathcal A(l_0)\qquad x\qquad x'
\]
\endgroup
